\section{Computer assistance and formal verification}\label{sec:lean}
\emph{What is computer-assisted.} The proof of Theorem~\ref{thm:hard} uses a computer at four finite steps (CA1)--(CA4) (Section~\ref{sec:hardcomp}): that every $P_a$ is a ZE partition of the thin chain for $7\le m\le60$, the transfer-matrix count $N_m=m$ for $7\le m\le40$ (the chain lemma for $m\ge10$ is computer-assisted, resting on a translation-invariance argument), the host profiles, and the parameter counts. Theorem~\ref{thm:sfA} uses a computer classification of $144$ configurations of binary relations, and Theorems~\ref{thm:CE} and~\ref{thm:MC0} use finite classifications; explicit small instances in the paper were checked by enumerating all set partitions. All other computations (independent re-implementations, random tests, exact optima of full thick instances for $k\le3$) are evidence, not part of any proof, and are listed in Appendix~\ref{app:checks}.

\emph{Lean~4.} Part of the argument was formalised in Lean~4 with Mathlib (project \texttt{P3Proofs}, directory \texttt{paper2\_p3\_interaction/\allowbreak lean/\allowbreak P3Proofs} of the repository \texttt{vinhqdang/\allowbreak correlation\_\allowbreak clustering} on GitHub). All files are free of \texttt{sorry} and of \texttt{native\_decide}; every theorem depends at most on the axioms \texttt{propext}, \texttt{Classical.choice} and \texttt{Quot.sound} (checked by \texttt{\#print axioms}); the finite computations use \texttt{decide +kernel}. The formalisation also found an error in a displayed inequality in the proof of Theorem~\ref{thm:hS} (corrected here).
\emph{Machine-checked:} the CSP representation of Section~\ref{sec:csp} (clusterings versus cluster graphs, $\mathrm{OPT}\ge|H|$, Observations~\ref{obs:one} and~\ref{obs:two} without the bound ``at most three owners'' and without the treewidth statements, and the five-state lemma); the twin lemma, half-degree and diameter at most $2$ of an optimal cluster and the separation of vertices at distance at least $3$, the lemma on the maximal family $M$ (MAX), the lemma on the hub (HUB, in the bounds $\nu(\Lambda_u)\le\ell$, vertex cover of size at most $2\ell$ and $|Y|\le2\ell+1$), and the small-merge lemma in abstract form with its hypotheses taken as assumptions; the block encoding of Grid Tiling and the equivalence of the comparator conditions with the Grid Tiling constraints (the last step of Proposition~\ref{prop:hTR}, $\Rightarrow$); the counting core of Theorem~\ref{thm:hS} as an abstract inequality system; the chain lemma for $m=7,8,9$ only; the algebra of the transfer matrix; and the parameter arithmetic of Theorem~\ref{thm:hard} including the no-split inequality.
\emph{Not machine-checked:} Theorem~\ref{thm:hS} at graph level; the blow-up lemma; host profiles (Lemma~\ref{lem:hprof}); tunnel release at graph level; the region bound (Lemma~\ref{lem:hregion}); pathwidth and pair-maximality of the reduction instance; the chain lemma for $m\ge10$; the link between the graph $\mathrm{OH}_m$ and the transfer matrix (so the steps (CA1)--(CA4) remain computer-assisted, with Lean covering only $m\le9$ and the matrix algebra); the dynamic programmes and running-time theorems; the ETH statements; and the reduction as a whole (Theorem~\ref{thm:hard}). The hypotheses of the counting core are taken for granted, so the Lean result certifies the arithmetic of Theorem~\ref{thm:hS}, not the claim that the instance $I_{\mathrm{tile}}$ satisfies them. A detailed list is in Appendix~\ref{app:checks}.
